\begin{definition}
	\textit{Линейным отображением} называется функция вида $w(z) = az + b$, где $a \in \Cm \bs \{0\}$, $b \in \Cm$, при этом считается, что $w(\infty) = \infty$.
\end{definition}

\begin{unnecessary}

\begin{note}
	Легко видеть, что линейное отображение непрерывно на $\CM$ и обратимо, причем обратное отображение $z(w) = \frac{w - b}{a}$ тоже линейно. Кроме того, композиция линейных отображений также является линейным отображением, поэтому линейные отображения с операцией композиции образуют группу.
\end{note}

\begin{example}
	Группа линейных отображений не является абелевой. Рассмотрим отображения $w = z + 1$ и $\zeta = 2z$, тогда $w(\zeta(z)) = 2z + 1$, но $\zeta(w(z)) = 2z + 2$.
\end{example}

\end{unnecessary}

\begin{note}
	Каждое линейное отображение конформно на $\CM$. Действительно, в конечных точках из $\Cm$ оно регулярно, а точка $\infty$ является его полюсом первого порядка.
\end{note}

\begin{definition}
	\textit{Дробно-линейным отображением} называется функция вида $w(z) \hm= \frac{az + b}{cz + d}$, где $a, b, c, d \in \Cm$, $ad - bc \ne 0$.
\end{definition}

\begin{note}
	В случае, когда в определении выше $c = 0$, получается уже линейное отображение, поэтому далее будем считать, что $c \ne 0$. Для такого дробно-линейного отображения можно также считать, что $w(\infty) = \frac ac$ и $w(-\frac dc) = \infty$.
\end{note}

\begin{proposition}
	Каждое дробно-линейное отображение обратимо, причем его обратное отображение тоже является дробно-линейным.
\end{proposition}

\begin{proof}
	Пусть $w(z) = \frac{az + b}{cz + d}$, тогда непосредственной подстановкой легко убедиться, что отображение $z(w) = \frac{dw - b}{-cw + a}$ является обратным к $w$.
\end{proof}

\begin{unnecessary}
    
\begin{corollary}
	Дробно-линейные отображения с операцией композиции образуют группу.
\end{corollary}

\begin{proof}
	Композиция дробно-линейных отображений также является дробно-линейным отображением, и каждое дробно-линейное отображение имеет обратное дробно-линейное отображение, что и означает требуемое.
\end{proof}

\begin{note}
	Назовем \textit{матрицей отображения} $w(z) = \frac{az + b}{cz + d}$ следующую матрицу:
	\[A_w := \begin{pmatrix} a&b\\c&d\end{pmatrix}\]
	
	Тогда при взятии композиции двух дробно-линейных отображения их матрицы перемножаются в соответствующем порядке.
\end{note}

\end{unnecessary}

\begin{proposition}
	Каждое дробно-линейное отображение конформно на $\CM$.
\end{proposition}

\begin{proof}
	Пусть $w(z) = \frac{az + b}{cz + d}$. Отображение $w$ однолистно на $\CM$ в силу обратимости, поэтому достаточно проверить, что оно конформно в каждой точке из $\CM$:
	\begin{itemize}
		\item Если $z \in \Cm \bs \{-\frac dc\}$, то $w$ регулярно в точке $z$, причём $w'(z) = \frac{ad - bc}{(cz + d)^2} \ne 0$
		
		\item Точка $z = -\frac dc$ является полюсом первого порядка отображения $w$
		
		\item Точка $z = \infty$ является устранимой особой точкой отображения $w$, причём выполнено условие $\res\limits_\infty w = \frac{ad - bc}{c^2} \ne 0$\qedhere
	\end{itemize}
\end{proof}

\begin{unnecessary}

\begin{proposition}
	Пусть $z_1, z_2, z_3 \in \CM$ и $w_1, w_2, w_3 \in \CM$  "--- тройки попарно различных точек. Тогда существует единственное дробно-линейное отображение $w$ такое, что $w(z_k) = w_k$ для каждого $k \in \{1, 2, 3\}$
\end{proposition}

\begin{proof}
	Сначала будем считать, что все точки из условия конечны. Рассмотрим отображения $\zeta(z) := \frac{z - z_1}{z - z_2}\cdot\frac{z_3 - z_2}{z_3 - z_1}$ и $\eta(z) := \frac{w - w_1}{w - w_2}\cdot\frac{w_3 - w_2}{w_3 - w_1}$. Тогда:
	\begin{center}
		\begin{tikzcd}[row sep = tiny, column sep = normal]
			z_1 \rar[maps to, "\zeta"] & 0 & \lar[maps to, swap, "\eta"]w_1
			\\
			z_2 \rar[maps to, "\zeta"] & \infty & \lar[maps to, swap, "\eta"]w_2
			\\
			z_3 \rar[maps to, "\zeta"] & 1 & \lar[maps to, swap, "\eta"]w_3
		\end{tikzcd}
	\end{center}
	
	Значит, отображение $w = \eta^{-1} \circ \zeta$ является искомым. Его удобнее задавать следующим равенством:
	\[\frac{z - z_1}{z - z_2}\cdot\frac{z_3 - z_2}{z_3 - z_1} = \frac{w - w_1}{w - w_2}\cdot\frac{w_3 - w_2}{w_3 - w_1}\]
	
	Если же среди точек из условия есть бесконечно удаленная точка, то формальное сокращение бесконечностей в равенстве выше дает новое равенство, задающее искомое отображение для этого случая. Опустим непосредственную проверку этого факта. \textit{Единственность лектор не доказал.}
\end{proof}

\end{unnecessary}

\begin{note}
	Образом прямой под действием линейного отображения всегда является прямая, а образом окружности под действием линейного отображения всегда является окружность. Далее мы получим аналогичное свойство для дробно-линейных отображений.
\end{note}

\begin{definition}
	\textit{Обобщённой окружностью} на $\CM$ называется множество, являющееся либо окружностью, либо объединением прямой и точки $\infty$.
\end{definition}

\begin{proposition}[круговое свойство дробно-линейных отображений]
	Образом обобщённой окружности под действием дробно-линейного отображения всегда является обобщённая окружность.
\end{proposition}

\begin{proof}
	Пусть $w(z) = \frac{az + b}{cz + d}$, тогда $w(z) = \frac ac - \frac{ad - bc}{c(cz + d)}$, то есть $w$ можно представить в виде композиции $w = z_3 \circ z_2 \circ z_1$, где $z_1(z) := cz + d$, $z_2(z_1) := \frac 1{z_1}$, $z_3(z_2) := \frac ac - \frac{ ad - bc}c z_2$. Значит, остается проверить, что отображение $z \mapsto \frac 1z$ переводит обобщённые окружности в обобщённые окружности. Известно, что любая прямая или окружность задается следующим уравнением с вещественными коэффициентами:
	\[A(x^2 + y^2) + Bx + Cy + D = 0,~B^2 + C^2 > 4AD\]
	
	Пусть $z = x + iy$, тогда $\frac 1z = \frac{x - iy}{x^2 + y^2} = \frac{x}{x^2 + y^2} - i\frac{y}{x^2 + y^2}$. Преобразование $(x, y) \mapsto \big(\frac{x}{x^2 + y^2}, -\frac{y}{x^2 + y^2}\big)$ даёт следующее уравнение:
	\[D(x^2 + y^2) + Bx - Cy + A = 0,~B^2 + C^2 > 4AD\]
	
	Как видно, это уравнение тоже задает прямую или окружность. Наконец, образом обобщённой окружности является вся обобщённая окружность, поскольку отображение $z \mapsto \frac 1z$ является обратным к самому себе.
\end{proof}